\documentclass[a4paper,12pt]{article}
\usepackage[utf8]{inputenc}
\usepackage[spanish]{babel}
\usepackage{amsmath}
\usepackage{geometry}

\geometry{margin=2.5cm}

\title{Solucionario: Examen de Física I - Forma B}
\date{}

\begin{document}
\maketitle

\section{Parte I: Teoría}
\begin{enumerate}
    \item \textbf{Etimología}: Del griego ``phisis'' (naturaleza/realidad).
    \item \textbf{Redondeo}: Menor cantidad de decimales en la operación.
    \item \textbf{Notación Científica}: Simplificar números extremadamente grandes o pequeños.
    \item \textbf{Magnitudes Derivadas}: Definidas por fundamentales. Ej: Velocidad (m/s), Fuerza (N).
\end{enumerate}

\section{Parte II: Práctica}
\begin{enumerate}
    \item \textbf{Caudal}: $45,000 / 3785.41 / 60 = 0.198$ gal/s.
    \item \textbf{Velocidad}: $2500 / 1.60934 \times 60 = 93,205.8$ mi/min.
    \item \textbf{Notación Científica}: $(8.4 / 2.1) \times 10^9 = 4.0 \times 10^9$.
    \item \textbf{Cifras Significativas}: $450.85 - 12.5 = 438.35 \rightarrow 438.4$.
\end{enumerate}

\end{document}
